% \begin{document}
\chapter{Restrizione e Induzione}

% ============================================================
\section{Rappresentazione Indotta}
% ============================================================

\subsection{Costruzione}

\dfn{Stabilizzatore di un Sottospazio}{
  Sia $ (V, \rho) $ una rappresentazione di $ G $ e $ W \subseteq V $ un sottospazio vettoriale. Lo \textit{stabilizzatore} di $ W $ è:
  \[
    H := \{g \in G \mid \rho(g)W \subseteq W\}
  \]
  Si verifica che $ H \leq G $ è un sottogruppo, e $ W $ è una rappresentazione di $ H $.
}

Costruiamo ora la \textbf{rappresentazione indotta} da $ W $.

Consideriamo la somma di tutti i traslati di $ W $ tramite $ G $:
\[
  U := \sum_{g \in G} \rho(g)W
\]

\nt{
  \textbf{Osservazione chiave}: $ \rho(g)W $ dipende solo dalla classe laterale $ gH $ di $ g $. Infatti se $ g' \in gH $, allora $ g' = gh $ per qualche $ h \in H $, e:
  \[
    \rho(g')W = \rho(gh)W = \rho(g)\rho(h)W = \rho(g)W
  \]
  perché $ \rho(h)W = W $ per definizione di $ H $. Quindi elementi dello stesso laterale danno lo stesso sottospazio, e la somma si può riscrivere sui laterali. Poiché laterali distinti danno sottospazi che si intersecano solo in $ \{0\} $, la somma è \textbf{diretta}:
  \[
    U = \bigoplus_{[g] \in G/H} \rho(g)W
  \]
}

\dfn{Rappresentazione Indotta}{
  $ U = \bigoplus_{[g] \in G/H} \rho(g)W $ si chiama \textbf{rappresentazione indotta} da $ W $, e si scrive:
  \[
    U = \mathrm{Ind}_H^G\, W
  \]
}

Si verifica che $ U $ è $ G $-stabile: per qualsiasi $ h \in G $ e qualsiasi addendo $ \rho(g)W $, si ha $ \rho(h)\cdot\rho(g)W = \rho(hg)W $, che è ancora un addendo della somma (quello corrispondente al laterale $ [hg] $).

\nt{
  La dimensione della rappresentazione indotta è:
  \[
    \dim U = |G/H| \cdot \dim W = \frac{|G|}{|H|} \cdot \dim W
  \]
}

\subsection{Applicazione: Bound sulla Dimensione delle Irriducibili}

\thm{Bound sulla Dimensione}{
  Se $ G $ è un gruppo finito e $ A \leq G $ è un sottogruppo abeliano, allora ogni rappresentazione irriducibile $ V $ di $ G $ soddisfa:
  \[
    \dim V \leq \frac{|G|}{|A|}
  \]
}

\pf{Dimostrazione}{
  \textbf{Passo 1.} Restringiamo l'azione a $ A $: $ V $ diventa una rappresentazione di $ A $.

  \textbf{Passo 2.} Per Maschke, $ V $ si decompone in irriducibili di $ A $. Poiché $ A $ è abeliano, per il Corollario di Schur ogni irriducibile di $ A $ ha dimensione 1 (ogni $ \rho(a) $ è un'omotetia, quindi ogni sottospazio 1-dimensionale è stabile). Quindi esiste $ W \leq V $ con $ \dim W = 1 $ che è un sotto-$ A $-modulo.

  \textbf{Passo 3.} Costruiamo la somma dei traslati:
  \[
    U = \sum_{[g] \in G/A} \rho(g)W \subseteq V
  \]
  La somma è indicizzata da $ G/A $ perché elementi dello stesso laterale $ gA $ danno lo stesso sottospazio ($ \rho(a)W = W $ per ogni $ a \in A $). Quindi $ U $ è un sottospazio $ G $-stabile generato da $ [G:A] $ sottospazi di dimensione 1, e quindi:
  \[
    \dim U \leq |G/A| \cdot \dim W = \frac{|G|}{|A|}
  \]

  \textbf{Passo 4.} $ U $ è un sotto-$ G $-modulo non nullo di $ V $. Poiché $ V $ è irriducibile, $ V = U $. Quindi:
  \[
    \dim V = \dim U \leq \frac{|G|}{|A|} \qquad\square
  \]
}

\cor{}{
  Ogni rappresentazione irriducibile di $ D_n $ ha dimensione $ \leq 2 $. \\
  }
\pf{Dimostrazione}{ $ D_n $ contiene il sottogruppo ciclico $ C_n $ di indice $ [D_n : C_n] = 2 $, quindi $ \dim V \leq |D_n|/|C_n| = 2 $
}

\subsection{Descrizione Esplicita tramite Classi Laterali}

Sia $R$ un sistema completo di rappresentanti per $G/H$ (un elemento di $G$ per ogni classe laterale). Ogni $g \in G$ permuta le classi laterali: dato $r \in R$,
\[
  gr = r_g h, \quad r_g \in R,\quad h \in H
\]
Quindi $g$ agisce su $V = \bigoplus_{r \in R} rW$ mandando $rW$ in $r_g W$: $r$ permuta gli addendi e $h$ agisce all'interno di $W$.

\nt{
  In modo più esplicito, $\forall r' \in R$, l'elemento $g$ manda $r'W$ in $gr'W = r_{g,r'}W$, dove $r_{g,r'}$ è il rappresentante di $g \cdot r'$ in $G/H$.
}

\ex{ES 1 — $H = \{e\}$, $W$ banale}{
  Se $H = \{e\}$, ci sono $|G|$ classi laterali (una per elemento). Quindi:
  \[
    \mathrm{Ind}_H^G W = \bigoplus_{g \in G} g \otimes \mathbb{C} \cong \mathbb{C}[G]
  \]
  Si ottiene la \textbf{rappresentazione regolare}! È come fare la rappresentazione regolare usando le classi laterali (invece di fissare tutti gli elementi di $G$).
}

\ex{ES 2 — $G = D_4$, $H = C_4 = \langle R \rangle$, $W = \mathbb{C}$}{
  Sia $W$ la rappresentazione di $C_4$ data da $R \cdot z = iz$ (quindi $R^k \cdot z = i^k z$).

  Le classi laterali sono $G/H = \{H,\, sH\}$, con rappresentanti $\{e, s\}$. Quindi:
  \[
    V = \mathrm{Ind}_H^G W = W \oplus sW
  \]
  con base $\{e_1,\, se_1\}$ dove $e_1$ è una base di $W$.

  Determiniamo l'azione di $R$ e $S$ (generatori di $D_4$) su $ae_1 + bse_1 \in V$.

  \textbf{Azione di $R$:} Usando la relazione $SR = R^{-1}S$ nel gruppo diedrale:
  \begin{align*}
    R \cdot (ae_1 + bse_1) &= aRe_1 + b(RS)e_1 \\
    &= a \cdot ie_1 + b \cdot SR^{-1}e_1 \\
    &= aie_1 + b(-i)se_1
  \end{align*}
  \[
    \Rightarrow\quad \rho(R) = \begin{pmatrix} i & 0 \\ 0 & -i \end{pmatrix}
  \]

  \textbf{Azione di $S$:}
  \begin{align*}
    S \cdot (ae_1 + bse_1) &= aSe_1 + bS^2e_1 = be_1 + ase_1
  \end{align*}
  \[
    \Rightarrow\quad \rho(S) = \begin{pmatrix} 0 & 1 \\ 1 & 0 \end{pmatrix}
  \]

  \nt{
    Poiché $\chi_U(R) = 0$, $\chi_U(R^2) = -2$, $\chi_U(S) = 0$, si verifica che $\chi_U = \chi_{V'}$ dove $V'$ è la rappresentazione naturale 2-dimensionale irriducibile di $D_4$. Quindi $U \cong V'$.
  }
}

% ============================================================
\subsection{Definizione tramite Prodotto Tensoriale}
% ============================================================

Prima di dare la definizione formale via tensori, ricordiamo la costruzione del prodotto tensoriale relativo.

\dfn{Prodotto Tensoriale di $A$-moduli}{
  Dati $V, W$ due $A$-moduli (su un anello $A$), il \textbf{prodotto tensoriale su $A$} è:
  \[
    V \otimes_A W \;:=\; \text{quoziente dell'$A$-modulo libero generato da } \{v \otimes w\}
  \]
  quozientato per le relazioni che rendono $\otimes$ $A$-bilineare (gli elementi di $A$ vengono trattati come scalari).
}

\ex{Esempi di prodotto tensoriale}{
  \begin{enumerate}
    \item $\mathbb{Z}^n \otimes_\mathbb{Z} \mathbb{Q} \cong \mathbb{Q}^n$: si prende il vettore intero e lo si ``immerge'' nei razionali.
    \item $(\mathbb{Z}, +) \otimes_\mathbb{Z} \mathbb{Q} = \{0\}$: sia $n = \mathrm{ord}(g)$, allora
      \[
        g \otimes \frac{p}{q} = ng \otimes \frac{p}{qn} = 0
      \]
      perché moltiplicare per $n$ dà $0$ a sinistra e dividere per $n$ non cambia il risultato a destra.
    \item $\mathbb{R}^n \otimes_\mathbb{R} \mathbb{C} \cong \mathbb{C}^n$: il tensore serve per \textbf{estendere gli scalari} da $\mathbb{R}$ a $\mathbb{C}$. Una base $\{x_1, \ldots, x_n\}$ di $\mathbb{R}^n$ diventa $\{x_1, \ldots, x_n, ix_1, \ldots, ix_n\}$.
  \end{enumerate}
}

\dfn{Rappresentazione Indotta (via tensore)}{
  \[
    \mathrm{Ind}_H^G W \;:=\; \mathbb{C}[G] \otimes_{\mathbb{C}[H]} W
  \]
  dove $\mathbb{C}[G]$ è visto come $\mathbb{C}[H]$-modulo a destra. L'azione di $G$ è:
  \[
    g \cdot (t \otimes w) = gt \otimes w, \qquad g \in G,\; t \in \mathbb{C}[G],\; w \in W
  \]
}

\nt{
  Se $g' \in gH$, cioè $g' = gh$, allora:
  \[
    g'(t \otimes w) = gh(t \otimes w) = ght \otimes w = gt \otimes hw
  \]
  Stiamo facendo il ``twisting'' in $H$. Questo è coerente: elementi della stessa classe laterale agiscono allo stesso modo (a meno dell'azione di $H$ su $W$). Si può vedere anche come un'estensione di scalari da $\mathbb{C}[H]$ a $\mathbb{C}[G]$.
}

\nt{
  La scelta del sistema di rappresentanti $R$ non è importante: la rappresentazione indotta è ben definita indipendentemente dalla scelta.
}

% ============================================================
\subsection{Equivalenza delle Due Definizioni}
% ============================================================

\thm{}{
  Sia $H \leq G$ e $(W, \sigma)$ una rappresentazione di $H$. Le due costruzioni della rappresentazione indotta sono isomorfe come $G$-moduli:
  \[
    \mathbb{C}[G] \otimes_{\mathbb{C}[H]} W \;\cong\; \bigoplus_{r \in R} rW
  \]
  dove $R$ è un sistema completo di rappresentanti per $G/H$.
}

\pf{Dimostrazione}{
  \textbf{Passo 1 — Decomposizione di $\mathbb{C}[G]$ come $\mathbb{C}[H]$-modulo destro.}

  Come $\mathbb{C}[H]$-modulo destro, $\mathbb{C}[G]$ si decompone secondo le classi laterali sinistre:
  \[
    \mathbb{C}[G] = \bigoplus_{r \in R} r \cdot \mathbb{C}[H]
  \]
  Infatti ogni $g \in G$ si scrive in modo unico come $g = rh$ con $r \in R$ e $h \in H$, quindi $\{rh \mid r \in R,\, h \in H\}$ è una base di $\mathbb{C}[G]$ e ogni $rh$ sta nell'addendo $r \cdot \mathbb{C}[H]$.

  \textbf{Passo 2 — Tensorizzazione addendo per addendo.}

  Il prodotto tensoriale commuta con le somme dirette (nel primo argomento), quindi:
  \[
    \mathbb{C}[G] \otimes_{\mathbb{C}[H]} W
    = \left(\bigoplus_{r \in R} r \cdot \mathbb{C}[H]\right) \otimes_{\mathbb{C}[H]} W
    \cong \bigoplus_{r \in R} \big(r \cdot \mathbb{C}[H] \otimes_{\mathbb{C}[H]} W\big)
  \]

  \textbf{Passo 3 — Ogni addendo è isomorfo a $W$.}

  Per ogni $r \in R$, consideriamo la mappa:
  \[
    \varphi_r : r \cdot \mathbb{C}[H] \otimes_{\mathbb{C}[H]} W \longrightarrow W, \qquad rh \otimes w \longmapsto \sigma(h)w
  \]
  Questa è ben definita: se spostiamo un elemento $h' \in H$ attraverso il tensore, abbiamo
  \[
    rh h' \otimes w = rh \otimes h'w = rh \otimes \sigma(h')w
  \]
  e la mappa dà $\sigma(hh')w = \sigma(h)\sigma(h')w$ nel primo caso e $\sigma(h)\sigma(h')w$ nel secondo, quindi è compatibile con le relazioni del tensore su $\mathbb{C}[H]$.

  La mappa inversa è $w \mapsto r \otimes w$, quindi $\varphi_r$ è un isomorfismo di spazi vettoriali. In particolare:
  \[
    \mathbb{C}[G] \otimes_{\mathbb{C}[H]} W \;\cong\; \bigoplus_{r \in R} W_r
  \]
  dove $W_r \cong W$ come spazio vettoriale, con l'elemento $r \otimes w$ che corrisponde a $w$ nella copia $W_r$.

  \textbf{Passo 4 — L'azione di $G$ coincide.}

  Nella costruzione tensoriale, $g \in G$ agisce per $g \cdot (r \otimes w) = gr \otimes w$. Scriviamo $gr = r' h_{g,r}$ con $r' \in R$ e $h_{g,r} \in H$. Allora:
  \[
    g \cdot (r \otimes w) = gr \otimes w = r' h_{g,r} \otimes w = r' \otimes \sigma(h_{g,r})w
  \]
  Questo corrisponde esattamente all'azione nella costruzione per somma diretta: $g$ manda l'addendo $W_r$ nell'addendo $W_{r'}$, applicando $\sigma(h_{g,r})$ all'interno di $W$. La mappa:
  \[
    \Phi : \mathbb{C}[G] \otimes_{\mathbb{C}[H]} W \longrightarrow \bigoplus_{r \in R} rW, \qquad r \otimes w \longmapsto \rho(r)w \in rW
  \]
  è un isomorfismo di $G$-moduli perché:
  \begin{align*}
    \Phi(g \cdot (r \otimes w)) &= \Phi(r' \otimes \sigma(h_{g,r})w) = \rho(r')\sigma(h_{g,r})w \\
    &= \rho(r' h_{g,r})w = \rho(gr)w = g \cdot \rho(r)w = g \cdot \Phi(r \otimes w) \qquad\square
  \end{align*}
}

% ============================================================
\subsection{Proprietà della Rappresentazione Indotta}
% ============================================================

\mprop{}{
  \begin{enumerate}
    \item $\dim \mathrm{Ind}_H^G W = \dim W \cdot \dfrac{|G|}{|H|}$
    \item Se $K \leq H \leq G$, allora $\mathrm{Ind}_H^G \mathrm{Ind}_K^H W \cong \mathrm{Ind}_K^G W$
  \end{enumerate}
}

\nt{
  Abbiamo quindi due modi per costruire rappresentazioni di $G$ a partire da sottogruppi:
  \begin{enumerate}
    \item \textbf{Discendere}: se $N \trianglelefteq G$, un $G/N$-modulo diventa $G$-modulo componendo con la proiezione $G \to G/N$. Es: $G = D_4$, $N = Z(G) \Rightarrow G/N \cong$ Klein.
    \item \textbf{Salire}: se $H \leq G$, un $H$-modulo diventa $G$-modulo tramite $\mathrm{Ind}_H^G$. Es: $G = D_4$, $H = C_4 \cong \mathbb{Z}/4$.
  \end{enumerate}
}

% ============================================================
\subsection{Carattere della Rappresentazione Indotta}
% ============================================================

\thm{Carattere dell'Indotta}{
  Sia $H \leq G$, $(W, \theta)$ rappresentazione di $H$ e $V = \mathrm{Ind}_H^G W$. Sia $R$ un sistema completo di rappresentanti per $G/H$. Allora per ogni $g \in G$:
  \[
    \chi_V(g) = \frac{1}{|H|} \sum_{\substack{s \in G \\ s^{-1}gs \in H}} \chi_W(s^{-1}gs)
  \]
}

\pf{Dimostrazione}{
  $V = \bigoplus_{r \in R} rW$ e $g$ permuta gli addendi: $\forall r \in R$, $gr = r_g h$ con $r_g \in R$, $h \in H$.

  Scegliamo una base per $V$ come unione di basi degli $rW$ al variare di $r \in R$ (possibile per somma diretta). Nella matrice di $\rho(g)$ in questa base:
  \begin{itemize}
    \item se $r_g \neq r$, l'addendo $rW$ viene mandato in un altro addendo $\Rightarrow$ nessun termine diagonale contribuisce;
    \item se $r_g = r$, cioè $gr = rh$ per qualche $h \in H$, allora $g$ manda $rW$ in sé stesso e la traccia del blocco è $\chi_W(h) = \chi_W(r^{-1}gr)$.
  \end{itemize}
  Quindi:
  \[
    \chi_V(g) = \sum_{\substack{r \in R \\ r_g = r}} \chi_W(r^{-1}gr)
      = \sum_{\substack{r \in R \\ r^{-1}gr \in H}} \chi_W(r^{-1}gr)
  \]
  Poiché questa somma dipende dalla rappresentazione $r$ delle classi, possiamo estenderla a tutti gli $s \in G$: all'interno della stessa classe laterale $rH$ ci sono $|H|$ elementi, tutti con lo stesso valore (il carattere è una funzione centrale), quindi dividendo per $|H|$:
  \[
    \chi_V(g) = \frac{1}{|H|} \sum_{\substack{s \in G \\ s^{-1}gs \in H}} \chi_W(s^{-1}gs) \qquad \square
  \]
}

% ============================================================
\section{Reciprocità di Frobenius}
% ============================================================

\subsection{Ripasso: Prodotto Hermitiano}

Ricordiamo che per $\varphi, \psi \in \mathcal{C}(G)$ si ha:
\[
  \langle \varphi, \psi \rangle_G = \frac{1}{|G|} \sum_{g \in G} \varphi(g)\,\overline{\psi(g)}
\]
Analogamente, per $\chi, \xi \in \mathcal{C}(H)$:
\[
  \langle \chi, \xi \rangle_H = \frac{1}{|H|} \sum_{h \in H} \chi(h)\,\overline{\xi(h)}
\]

\subsection{Restrizione e Induzione}

Sia $H \leq G$. Abbiamo due operazioni:
\begin{itemize}
  \item \textbf{Induzione}: da una rappresentazione $(W, \theta)$ di $H$ costruiamo $\mathrm{Ind}_H^G W$ (rappresentazione di $G$), con carattere $\mathrm{Ind}\,\chi$.
  \item \textbf{Restrizione}: da una rappresentazione $(V, \rho)$ di $G$ otteniamo $\mathrm{Res}_H^G V$ (rappresentazione di $H$), componendo $H \hookrightarrow G \xrightarrow{\rho} GL(V)$, con carattere $\mathrm{Res}_H^G \rho$.
\end{itemize}
Possiamo quindi sia restringere $G$ ad $H$, sia far ``salire'' $H$ a $G$. Il teorema di Frobenius dice che è la stessa cosa.

\thm{Reciprocità di Frobenius}{
  Siano $H \leq G$ un sottogruppo di un gruppo finito $G$, $W$ una rappresentazione di $H$ e $U$ una rappresentazione di $G$. Allora:
  \[
    \langle \chi_{\mathrm{Ind}_H^G W},\; \chi_U \rangle_G \;=\; \langle \chi_W,\; \chi_{\mathrm{Res}_H^G U} \rangle_H
  \]
}

\pf{Dimostrazione}{
  Calcoliamo il lato sinistro usando il teorema sul carattere dell'indotta:
  \begin{align*}
    \langle \chi_{\mathrm{Ind}_H^G W}, \chi_U \rangle_G
    &= \frac{1}{|G|} \sum_{g \in G} \chi_{\mathrm{Ind}_H^G W}(g)\,\overline{\chi_U(g)} \\
    &\overset{\text{Thm}}{=} \frac{1}{|G|} \cdot \frac{1}{|H|} \sum_{g \in G} \sum_{\substack{s \in G \\ s^{-1}gs \in H}} \chi_W(s^{-1}gs)\,\overline{\chi_U(g)}
  \end{align*}
  Facciamo il cambio di variabili $s^{-1}gs = h \in H$, cioè $g = shs^{-1}$. Poiché $\chi_U$ è una funzione centrale, $\overline{\chi_U(g)} = \overline{\chi_U(shs^{-1})} = \overline{\chi_U(h)}$. La somma diventa:
  \[
    = \frac{1}{|G|\,|H|} \sum_{s \in G} \sum_{h \in H} \chi_W(h)\,\overline{\chi_U(h)}
  \]
  Poiché $s$ non compare nell'integrando, la somma su $s$ dà un fattore $|G|$:
  \[
    = \frac{1}{|H|} \sum_{h \in H} \chi_W(h)\,\overline{\chi_U(h)} = \langle \chi_W,\, \chi_{\mathrm{Res}_H^G U} \rangle_H \qquad \square
  \]
}

\ex{Esempio: $G = D_4$, $V$ naturale, $W$ repr. di $C_4$}{
  Sia $V$ la rappresentazione naturale di $D_4$ su $\mathbb{C}^2$ e sia $W$ la rappresentazione di $C_4 = \langle R \rangle \leq D_4$ data da $R \cdot z = iz$.

  Vogliamo calcolare $\langle \chi_V, \mathrm{Ind}\,\chi_W \rangle_{D_4}$ (che vale $1 \iff V \cong \mathrm{Ind}\,\chi_W$).

  Per Reciprocità di Frobenius:
  \[
    \langle \chi_V, \mathrm{Ind}\,\chi_W \rangle_{D_4}
    = \langle \mathrm{Res}_{C_4}^{D_4}\,\chi_V,\, \chi_W \rangle_{C_4}
    = \frac{1}{4} \sum_{R^i \in C_4} \chi_V(R^i)\,\overline{\chi_W(R^i)}
  \]
  I valori del carattere di $V$ sulle rotazioni (dalla tavola di $D_4$) sono $\chi_V(e) = 2$, $\chi_V(R) = 0$, $\chi_V(R^2) = -2$, $\chi_V(R^3) = 0$. I valori di $\chi_W$ sono $\chi_W(R^k) = i^k$. Quindi:
  \[
    = \frac{1}{4}\left(2 \cdot 1 + 0 \cdot (-i) + (-2) \cdot (-1) + 0 \cdot i\right) = \frac{1}{4}(2 + 0 + 2 + 0) = 1
  \]
  Quindi $V \cong \mathrm{Ind}_{C_4}^{D_4} W$. \checkmark
}
